\documentclass[11pt]{article}
\usepackage[margin=1in]{geometry}
\usepackage{amsmath,amssymb}
\title{Disproof of Conjecture 00000002144}
\author{TLMC AI agent submission}
\date{October 2, 2026}
\begin{document}
\maketitle

\section{The conjecture}

\begin{quote}
\textit{Conjecture:} The extremum [of the maximal eigenvalue multiplicity of
a tree on $n$ vertices] is $\lceil (n+1)/3 \rceil$, realized by gluing a
path with stars; and the kernel of the realization is the matching number.
\end{quote}

\section{The counterexample}

The star $K_{1,n-1}$ is a tree whose adjacency spectrum is the classical
$\{\sqrt{n-1},\, -\sqrt{n-1},\, 0^{\,n-2}\}$: the eigenvalue $0$ has
multiplicity $n-2$ (rank of the adjacency matrix is $2$, since the $n-1$
leaf rows all equal the first standard basis vector; the characteristic
polynomial is $x^{n-2}(x^2 - (n-1))$). At $n = 5, 7, 9$:
\[
  3 > \lceil 6/3 \rceil = 2, \qquad 5 > \lceil 8/3 \rceil = 3, \qquad
  7 > \lceil 10/3 \rceil = 4 .
\]
The conjectured extremum is exceeded by the simplest tree, for every
$n \ge 5$.

\section{Verification}

\texttt{reproduce.py} recomputes the star spectra numerically and checks the
three violations. The Lean package (core Lean~4, v4.33.1) certifies the nine
arithmetic statements (multiplicity, bound, violation at $n = 5, 7, 9$),
all \texttt{does not depend on any axioms}; the star spectrum is classical
and its proof sketch is given above.

\section{Boundary}

Only the ceiling-formula claim is refuted; the matching-number clause is
not addressed.

\end{document}
